%=======================================================================
%  CSE558 Data Science - Midterm Project Report (built on the official template)
%  Monsoon 2026, IIIT-Delhi.  Group DSC 10.
%  Every number below is produced by 1_midterm/analysis/*.py; see
%  1_midterm/results/audit_numbers.csv for number -> source file.
%  Compile with pdfLaTeX, twice.
%=======================================================================
\documentclass[10pt,a4paper,twocolumn]{article}

%========================= EDIT THIS =========================
\newcommand{\projecttitle}{Promotion timing, missing data and competitors:
  four hypothesis tests on 1{,}115 retail hubs}
\newif\ifshowhints
\showhintsfalse      % \showhintsfalse hides every grey instruction line
%=============================================================

%-------------------------------------------- fonts (no need to edit)
\usepackage{iftex}
\ifPDFTeX
  \usepackage[T1]{fontenc}
  \usepackage[utf8]{inputenc}
  \IfFileExists{carlito.sty}%
    {\usepackage[sfdefault]{carlito}}%              Calibri look
    {\usepackage[scaled=0.92]{helvet}\renewcommand{\familydefault}{\sfdefault}}%
  \usepackage{courier}
\else
  \usepackage[no-math]{fontspec}
  \setmainfont{Carlito}[Extension=.ttf, UprightFont=*-Regular, BoldFont=*-Bold,
    ItalicFont=*-Italic, BoldItalicFont=*-BoldItalic]
\fi

%----------------------------------------------------------- packages
\usepackage[margin=18mm,bottom=20mm,columnsep=7mm]{geometry}
\usepackage{amsmath}
\usepackage{xcolor}
\usepackage{graphicx}
\usepackage{tabularx}
\usepackage{array}
\usepackage{booktabs}
\usepackage{enumitem}
\usepackage{titlesec}
\usepackage{needspace}
\usepackage[font=small,labelfont=bf]{caption}
\usepackage[hidelinks]{hyperref}
\hypersetup{pdftitle={\projecttitle},
            pdfsubject={CSE558 Data Science, Monsoon 2026: Midterm Project Report}}

%------------------------------------------------------------- layout
\definecolor{eg}{HTML}{7396C4}         % worked example (light blue)
\definecolor{hintgrey}{gray}{0.58}     % instructions (light grey)
\definecolor{keyblue}{HTML}{0B5FA5}   % key values and headings (colour-blind safe pair)
\definecolor{warnor}{HTML}{C2410C}    % caveats and fragile results

\titleformat{\section}{\normalfont\large\bfseries\color{keyblue}}{\thesection.}{0.5em}{}
\titlespacing*{\section}{0pt}{12pt}{3pt}
\titleformat{\subsection}{\normalfont\normalsize\bfseries\color{keyblue}}{\thesubsection}{0.5em}{}
\titlespacing*{\subsection}{0pt}{10pt}{2pt}
\newcommand{\sectionbreak}{\needspace{10\baselineskip}}
\newcommand{\subsectionbreak}{\needspace{10\baselineskip}}

\setlength{\parindent}{0pt}
\setlength{\parskip}{4pt}
\setlist{leftmargin=1.2em, itemsep=2pt, topsep=2pt, parsep=0pt}
\renewcommand{\arraystretch}{1.25}
\setlength{\tabcolsep}{4pt}
\newcolumntype{L}[1]{>{\raggedright\arraybackslash}p{#1}}
\newcolumntype{Y}{>{\raggedright\arraybackslash}X}
\newcolumntype{K}[1]{>{\raggedright\arraybackslash\bfseries\boldmath}p{#1}}  % bold labels

%--------------------------------------------------------- shorthands
\newcommand{\hint}[1]{\ifshowhints{\footnotesize\itshape\color{hintgrey}#1\par}\fi}
\newcommand{\key}[1]{\textbf{\textcolor{keyblue}{#1}}}
\newcommand{\warn}[1]{\textbf{\textcolor{warnor}{#1}}}
\newcommand{\field}[2]{\key{#1}~#2\par}
\newcommand{\tpart}[2]{\par\needspace{12\baselineskip}\smallskip\key{#1)~#2}\par}
\newcommand{\minihead}[1]{\par\needspace{8\baselineskip}\smallskip\key{#1}\par}
\newcommand{\code}[1]{\texttt{#1}}
\newcommand{\figfile}[2]{%
  \IfFileExists{#1}%
    {\includegraphics[width=\linewidth,height=#2,keepaspectratio]{#1}}%
    {\fbox{\parbox[c][#2][c]{\dimexpr\linewidth-2\fboxsep-2\fboxrule\relax}%
      {\centering\small\itshape\color{hintgrey}Your plot: \detokenize{#1}}}}}
\newcommand{\FIG}{../figures/}

%=======================================================================
\begin{document}

\twocolumn[%
{\small\bfseries CSE558 Data Science\ \ \textbullet\ \ Midterm Project Report\ \ \textbullet\ \ Monsoon 2026\par}
\vspace{4pt}
{\LARGE\bfseries\raggedright\projecttitle\par}
\vspace{6pt}
\rule{\linewidth}{1pt}

{\small
\begin{tabularx}{\linewidth}{@{}K{2.4cm} Y@{}}
\toprule
Group name & DSC 10 \\
Dataset & Shadow Hub database (training, test and hub files, 1{,}115 hubs):
  \url{https://github.com/kushiiitd05/ShadowDsc10-ShadowHub/tree/cse558-midterm-2026/1_midterm/data/raw} \\
Code & \url{https://github.com/kushiiitd05/ShadowDsc10-ShadowHub/tree/cse558-midterm-2026}
  (scripts in \code{1\_midterm/analysis/}, README in \code{1\_midterm/README.md}) \\
\bottomrule
\end{tabularx}}

\vspace{2pt}
{\small
\begin{tabularx}{\linewidth}{@{}L{0.5cm} L{3.4cm} L{2.3cm} L{4.8cm} Y@{}}
\toprule
\textbf{\#} & \textbf{Name} & \textbf{Roll no.} & \textbf{Institute email} & \textbf{Contribution} \\
\midrule
1 & Manya Chaturvedi & 2023312 & manya23312@iiitd.ac.in & <<CONFIRM: Dataset and size check (1)>> \\
2 & Kush Tokas & 2023296 & kush23296@iiitd.ac.in & <<CONFIRM: Preparation, before and after (2, 3)>> \\
3 & Akash Adur & 2024053 & akash24053@iiitd.ac.in & <<CONFIRM: EDA and problem statements (4)>> \\
4 & Anusha Rana & 2023114 & anusha23114@iiitd.ac.in & <<CONFIRM: Tests 1 and 2 (5)>> \\
5 & Sanjeev & 2023483 & sanjeev23483@iiitd.ac.in & <<CONFIRM: Tests 3 and 4, limitations, code and README (5, 6)>> \\
\bottomrule
\end{tabularx}}
\vspace{10pt}
]

%-----------------------------------------------------------------------
\section{Dataset}

\field{Source:}{the training, test and hub files of the Shadow Hub database, which our group received through a private
data-science competition; the files are in our repository (link above). One row of \code{orders\_train.csv} is one retail hub
on one calendar day, 1~January~2013 to 19~June~2015 (900 dates, 1{,}115 hubs). \code{hub\_metadata.csv} has one row per hub.
\code{orders\_test.csv} (20~June to 31~July~2015) has no outcome, so it is used only to describe the schedule.}

\field{Provenance:}{the files match the public Rossmann Store Sales data with renamed columns. Evidence:
1{,}115 hubs; $970{,}379+46{,}830=1{,}017{,}209$ rows, Rossmann's training size; hub~1 (format~3, tier~1,
competitor distance 1{,}270, competitor since September 2008, no loyalty programme) and hub~2 (loyalty since week 13
of 2010, months ``Jan,Apr,Jul,Oct'') reproduce Rossmann stores~1 and~2; and the same 180 hubs are missing from
1~July to 31~December~2014.}

{\small
\begin{tabularx}{\linewidth}{@{}L{3.0cm} Y Y@{}}
\toprule
 & \textbf{Raw file} & \textbf{After encoding} \\
\midrule
Instances (rows) & 970{,}379 & 970{,}379 \\
Attributes (columns) & 9 (+ 9 hub fields) & 46 \\
Instances $\times$ attributes & 8{,}733{,}411 & \key{44{,}637{,}434} \\
\bottomrule
\end{tabularx}}

\field{Size check:}{\key{$970{,}379\times46=44{,}637{,}434\ge3{,}000{,}000$}; the raw file alone gives
$970{,}379\times9=8{,}733{,}411$. The 46 attributes are: \code{HubID}; 4 calendar fields from \code{Date};
one-hot \code{Weekday} (7), \code{RegionalHoliday} (4), \code{HubFormat} (4) and loyalty interval (4, incl.\ ``None'');
14 numeric or binary fields; and 8 data-quality flags (Section 2). Rows are counted as hub-days, but every test uses
the \emph{hub} as its independent unit (Section 5).}

\minihead{Key features}
{\footnotesize
\begin{tabularx}{\linewidth}{@{}L{3.25cm} L{1.35cm} Y@{}}
\toprule
\textbf{Feature} & \textbf{Type} & \textbf{Real-world meaning (verified facts)} \\
\midrule
\code{OrderVolume} & discrete & \key{response}: the hub's daily volume, \emph{units unstated}. Median volume per app
  session is 9.23 on open days, which is implausible for a count of orders \\
\code{AppSessions} & discrete & app sessions (traffic); Pearson $r=0.825$ with volume on open days \\
\code{IsOpen} & nominal & hub open that day (82.87\% of rows) \\
\code{PromoActive} & nominal & promotion running (38.27\%); set per date, identical for all hubs, in alternate
  Monday--Friday weeks, never at weekends \\
\code{Weekday} & ordinal (cyclic) & 1 = Monday \dots\ 7 = Sunday; matches \code{Date} in every row \\
\code{RegionalHoliday} & nominal (4) & none, public holiday, Easter, Christmas \\
\code{HubFormat} & nominal (4) & store format; 602 / 17 / 148 / 348 hubs \\
\code{CompetitorDistance} & continuous & distance to the nearest competitor (median 2{,}325; 3 hubs missing) \\
\code{CompetitorOpenSince}\newline\code{Month, Year} & discrete & month and year that competitor opened; unknown for 354 hubs \\
\bottomrule
\end{tabularx}}

%-----------------------------------------------------------------------
\section{Data preparation}

No raw value was overwritten; every step adds a flag or a new column, so each choice can be undone.
``Why not'' names the alternative we rejected.

{\scriptsize
\renewcommand{\arraystretch}{1.12}
\begin{tabularx}{\linewidth}{@{}L{2.05cm} L{2.0cm} Y@{}}
\toprule
\textbf{Issue found} & \textbf{What we did} & \textbf{Why, and not the alternative} \\
\midrule
\key{33{,}121} absent hub-dates: 180 hubs $\times$ 184 days (2014-H2) + hub 988 on 2013-01-01
 & left absent; flag \code{gap\_hub}; tested in T2
 & a systematic block, not random gaps. Zero-filling would invent 33{,}121 closures; mean, median, kNN or regression
   imputation would invent 2014-H2 sales for exactly the hubs that T2 shows are not a random subset \\
\key{166{,}269} closed rows, all with volume 0 and sessions 0
 & kept, excluded from volume tests
 & structural zeros, neither missing nor outliers; mixing ``is it open?'' with ``how much?'' gives a two-humped mixture \\
54 open rows with volume 0 (41 hubs)
 & flag \code{open\_zero}; excluded; listed in README
 & an open day with no activity is inconsistent. MLE $\hat p=54/804{,}110=6.7\times10^{-5}$ (Wilson 95\% CI
   $5.1\times10^{-5}$ to $8.8\times10^{-5}$), so dropping them cannot move any estimate \\
Pooled IQR fence (Q1 4{,}856, Q3 8{,}360, fence 13{,}616) flags \key{29{,}307} rows (3.64\%)
 & not deleted; tests use $\log(1+V)$ and per-hub contrasts
 & \key{88.3\%} of these rows belong to the 189 hubs whose \emph{own} median exceeds Q3: the fence measures hub size,
   not error. Trimming or winsorising would cut big hubs' ordinary days \\
Per-hub robust $z$ (median/MAD, $|z|>3.5$): 6{,}220 rows, 1{,}040 hubs
 & kept as genuine demand
 & 97.0\% sit next to a closure (63.2\%), before Christmas (17.5\%), on or next to a holiday (11.1\%), on promo
   (4.2\%) or after a reopening (1.1\%). The maximum, 38{,}722 (hub 262), is Good Friday 2015 with a promotion \\
Normal ``3$\sigma$'' rule
 & not used
 & $|z|\ge3$ holds for 1.59\% of open days against 0.27\% under normality (kurtosis 7.84); the distribution-free
   Chebyshev bound (11.1\%) holds, but a normal cut-off would mislabel real peaks \\
88 closure runs $\ge5$ days at 85 hubs (longest 185)
 & flags \code{long\_closure}, 14-day \code{reopen\_window}
 & refurbishment-like; volume in the 14 days after reopening is 44\% above the 28 days before closing, so reopen days
   are excluded in a sensitivity rerun \\
Volume/session ratio: median 9.23, max 64.96 (hub 897: 16{,}954 / 261); 23 rows $>30$
 & flag \code{ratio\_anomaly}; kept
 & sessions are probably under-recorded; no test uses sessions \\
Metadata: distance missing (3 hubs); opening date missing (354); year 1900 (hub 815); 1961 (hub 146)
 & complete-case in T3 ($n=1{,}112$); date left NA + indicator; 1900 $\to$ NA; 1961 kept
 & imputation would place hubs at invented distances or dates inside the very tests that use them; a 1900 opening
   means a 115-year-old competitor, a sentinel \\
Loyalty start/interval NA for 544 hubs
 & level ``None''
 & exactly the 544 hubs without the programme: \emph{not applicable}, not missing \\
Duplicates
 & none removed
 & 0 duplicate hub-dates, 0 duplicate rows; 8 hubs form 4 pairs with identical metadata but different daily series, so
   they are distinct hubs \\
Integer codes for nominal fields
 & one-hot; tier as ordinal code; log scale as a new column
 & codes 1--4 for format would impose a false order and distance \\
Integrity
 & checked
 & last digit of volume ($\ge100$) is uniform, $\chi^2(9)=11.10$, $p=0.27$: no heaping or rounding \\
\bottomrule
\end{tabularx}}

%-----------------------------------------------------------------------
\section{What changed after preparation}

{\footnotesize
\begin{tabularx}{\linewidth}{@{}L{3.0cm} L{1.55cm} Y@{}}
\toprule
\textbf{Measure} & \textbf{Before} & \textbf{After} \\
\midrule
Rows / attributes & 970{,}379 / 9 & 970{,}379 / 46 (no row deleted) \\
Rows that enter the tests & -- & 804{,}056 open rows with volume $>0$ \\
Missing cells after the hub join & 2{,}075{,}025 & 1{,}591{,}274, of which 970{,}366 are \emph{not applicable} loyalty fields \\
Skewness of open-day volume & 1.59 (raw) & $-0.11$ on $\log(1+V)$ \\
Kurtosis (normal $=3$) & 7.84 (raw) & 3.65 on $\log(1+V)$ \\
Flagged rows & -- & open-zero 54; ratio 23; long closure 1{,}898; reopen window 1{,}218 \\
Second outlier pass (robust $z$ on $\log(1+V)$) & -- & 2{,}645 rows (2{,}060 low, 585 high); 1{,}870 explained by holiday, closure, reopening or pre-Christmas; the rest kept \\
\bottomrule
\end{tabularx}}

\key{Nothing was deleted:} preparation turned silent problems into flags, and each test states which flags it excludes. The
log scale removes most of the skew and the excess kurtosis, which is why every test uses it. The second outlier pass,
run after the fixes as Lecture~3 asks, mostly flags \emph{low} days next to closures and holidays, so it does not justify
any further deletion.

%-----------------------------------------------------------------------
\section{Basic questions and problem statement}

\begin{figure}[!ht]
  \figfile{\FIG fig_volume.pdf}{3.6cm}
  \caption{Open-day volume is right-skewed (skewness 1.59, kurtosis 7.84); on $\log(1+V)$ it is nearly symmetric
  ($-0.11$, 3.65), \key{so all tests use the log scale}.}
\end{figure}

\begin{figure}[!ht]
  \figfile{\FIG fig_gap.pdf}{3.1cm}
  \caption{\key{180 of 1{,}115 hubs (16.1\%)} have no rows at all in 2014-H2; \key{format~4 loses 23.3\%} of its hubs (Wilson 95\%
  CIs; hubs per format in brackets).}
\end{figure}

\begin{figure}[!ht]
  \figfile{\FIG fig_promo_weekday.pdf}{2.9cm}
  \caption{The promotion lift \key{falls through the week, from 60\% on Monday to 22\% on Friday} (mean of 1{,}115 hub-level
  lifts; 95\% CIs are shorter than the dots).}
\end{figure}

\begin{figure}[!ht]
  \figfile{\FIG fig_distance.pdf}{3.4cm}
  \caption{Hub volume is flat in competitor distance except in the closest decile (distance $\le250$), whose 122 hubs
  \key{sell 15.9\% more} than the rest.}
\end{figure}

\minihead{Basic questions}
\begin{enumerate}[label=\arabic*.]
  \item Does a promotion raise volume? (Yes: by 22--60\%, depending on the weekday, Figure~3.)
  \item Which format sells most? (Format~2: its 17 hubs average 10{,}078 per open day; formats 1, 3, 4 average 6{,}816--6{,}919.) Are Sundays different? (Only 33 hubs ever open on Sunday.)
  \item Do volume and app sessions move together? ($r=0.825$ on open days.)
\end{enumerate}

\minihead{Our problem statements}
\noindent\fbox{\begin{minipage}{\dimexpr\linewidth-2\fboxsep-2\fboxrule\relax}
\key{P1 (primary).} Within the same hub, is the promotion lift on Mondays different from the lift on Fridays?

\smallskip
\key{P2.} Is the 184-day data gap independent of hub format, i.e.\ is the missingness random?

\smallskip
\key{P3.} Is the distance to the nearest competitor associated with a hub's typical volume?

\smallskip
\key{P4.} When a competitor opens, does the hub's volume change relative to hubs with no new competitor over the
same weeks?
\end{minipage}}

\field{How they are novel:}{the basic questions read one row or one column at a time. Ours need a contrast that no
row records: a difference of lifts inside each hub (P1), the pattern of \emph{missing} rows (P2), and a hub's change
measured against control hubs over the same calendar weeks (P4). P3 aggregates two years to one value per hub.}

\minihead{Train vs test schedule (the test file has no outcome)}
{\footnotesize
\begin{tabularx}{\linewidth}{@{}Y L{1.0cm} L{1.0cm} L{1.25cm} L{1.0cm}@{}}
\toprule
\textbf{Rows} & \textbf{Open} & \textbf{Promo} & \textbf{School closed} & \textbf{Holiday} \\
\midrule
Train, all 900 days & 82.9\% & 38.3\% & 17.4\% & 3.2\% \\
Train, last 42 days & 80.0\% & 35.7\% & 7.5\% & 6.4\% \\
Train, 20 Jun--31 Jul 2014 & 85.6\% & 34.9\% & 27.6\% & 0.0\% \\
Test, 20 Jun--31 Jul 2015 & 86.0\% & 35.7\% & 28.5\% & 0.0\% \\
\bottomrule
\end{tabularx}}

\field{Two facts that shape the tests.}{(i) Days of one hub are dependent: the median within-hub correlation of
$\log(1+V)$ is 0.65 at lag~1 and 0.55 at lag~14, but only 0.02 at lag~7, because promotions alternate weekly.
(ii) The test window is the summer school break (school closure 28.5\% vs 17.4\% overall; 27.6\% in the same weeks of
2014), so it is seasonal rather than a new regime.}

%-----------------------------------------------------------------------
\section{Statistical inferences}

\field{Common to all tests.}{$y=\ln(1+\text{OrderVolume})$ on open days with volume $>0$. Because hub-days are
dependent (Section~4), each test reduces the data to \emph{one number per hub} ($n$ hubs). We write $\bar x$ for a
sample mean, $s$ for a sample standard deviation (divisor $n-1$), $SE$ for a standard error, $\mu$, $\rho$, $\pi$ for
population values, and $\delta$ for the smallest difference that matters. Every statistic is $(\text{estimate}-\text{null
value})/SE$, except $\chi^2$. $\alpha=0.05$. All tests are \textbf{two-sided}: an earlier version of these tests on raw
volume had already shown the directions, so choosing a side now would be choosing it after seeing the data. We fixed the
analysis choices in a committed plan (\code{analysis\_plan\_r1.md}) before computing these statistics. With four tests,
Bonferroni uses $\alpha/4=0.0125$. Every statistic was computed with \code{scipy.stats} and by the formula shown;
they agree to 4 significant figures. The $\delta$ values are our assumptions, with the reason given.}

%------------------------------------------------------------ Test 1
\subsection*{Test 1 (P1): is the promotion lift different on Monday and Friday?}

\tpart{a}{Hypotheses, written before running the test}
{\small
\begin{tabularx}{\linewidth}{@{}K{3.4cm} Y@{}}
\toprule
$H_0$ & $\mu_d=0$ \\
$H_1$ & $\mu_d\neq0$ \\
Significance level / side & $\alpha=0.05$, two-sided \\
Test & paired $t$-test (one-sample $t$ on $d_i$) \\
Smallest difference that matters & $\delta=0.05$ on the log scale ($\approx5\%$ of volume) \\
\bottomrule
\end{tabularx}}

For hub $i$, $L_{i,\text{Mon}}=\bar y_{i,\text{Mon,promo}}-\bar y_{i,\text{Mon,no promo}}$ and likewise $L_{i,\text{Fri}}$,
on open, non-holiday days outside reopening windows (each of the 4 cells $\ge4$ days); $d_i=L_{i,\text{Mon}}-L_{i,\text{Fri}}$
and $\mu_d$ is its population mean. $e^{d}$ is the ratio of the Monday promotion multiplier ($e^{L_{i,\text{Mon}}}$) to the Friday one, so $e^{d}-1$ is how much larger the Monday gain is. $\delta=5\%$:
below that, the choice of which weekday carries the promotion budget would not change.

\tpart{b}{Other tests we considered}
The same hub appears in all four cells, so the data are paired; $\sigma$ is unknown, so we use $t$.
{\footnotesize
\begin{tabularx}{\linewidth}{@{}L{2.4cm} Y@{}}
\toprule
\textbf{Also considered} & \textbf{Why we did not use it} \\
\midrule
Welch $t$ on hub-days & treats 262{,}170 dependent days (lag-1 $r=0.65$) as independent; $p$ would be far too small \\
Pooled two-sample $t$ & the two weekdays come from the same hubs, so the groups are not independent \\
$z$-test & needs the population $\sigma$ \\
Wilcoxon signed-rank, sign test & kept as robustness checks: $p=6.0\times10^{-184}$; $d_i>0$ for 1{,}112 of 1{,}115 hubs \\
\bottomrule
\end{tabularx}}

\tpart{c}{Result}
\[t=\frac{\bar d-0}{s_d/\sqrt n}=\frac{0.2691}{0.09932/\sqrt{1115}}=\frac{0.2691}{0.002974}=90.47.\]
{\small
\begin{tabularx}{\linewidth}{@{}K{3.4cm} Y@{}}
\toprule
Estimate & $\bar d=0.269$: \key{Monday's promotion gain is 30.9\% larger} (volume $\times1.600$ on Monday, $\times1.222$ on Friday; ratio $1.600/1.222=1.309$) \\
95\% confidence interval & $0.263$ to $0.275$ (30.1\% to 31.6\%) \\
Test statistic (d.f.) & $t=90.47$ (1{,}114); $t_{0.975,1114}=1.962$ \\
p-value & $<10^{-300}$ \\
Decision & \key{reject $H_0$} (also under Bonferroni) \\
\bottomrule
\end{tabularx}}

The whole interval lies beyond $\delta=0.05$, so the difference is large enough to matter. The normal Q-Q plot of $d_i$
is close to a straight line (Figure~5), and with $n=1{,}115$ the CLT makes $\bar d$ close to normal.
\textbf{Replication:} 2013 only, $\bar d=0.249$; 2014--15 only, $0.286$; odd and even HubID, $0.267$ and $0.272$; with
holidays and reopening days included, $0.270$. All reject. On raw volume, $+2{,}169$ (95\% CI $2{,}105$--$2{,}233$).
\textbf{Placebo:} shuffling the promotion labels within hub $\times$ weekday 200 times gives a mean estimate of 0.0002;
7.0\% of the 200 shuffles reject at 0.05, close to the 5\% expected.
\textbf{Shared calendar check:} the promotion calendar is the same for every hub, so hubs share date shocks. Taking the
\emph{week} as the unit (64 promotion and 52 non-promotion weeks, contrast = cross-hub Monday mean minus Friday
mean): an $F$-test rejects equal variances ($F=0.435$, d.f.\ 63 and 51, $p=0.0018$), so we use Welch: $t=10.6$
(d.f.\ 84.9), estimate 0.284, 95\% CI $0.231$--$0.337$, $p=2.7\times10^{-17}$. The interval is 9 times wider, but it
still lies wholly beyond $\delta$.

%------------------------------------------------------------ Test 2
\subsection*{Test 2 (P2): is the 2014-H2 data gap independent of hub format?}

\tpart{a}{Hypotheses, written before running the test}
{\small
\begin{tabularx}{\linewidth}{@{}K{3.4cm} Y@{}}
\toprule
$H_0$ & $\pi_1=\pi_2=\pi_3=\pi_4$ (gap independent of format) \\
$H_1$ & not all $\pi_f$ are equal \\
Significance level / side & $\alpha=0.05$; non-directional (upper tail of $\chi^2$) \\
Test & $\chi^2$ test of independence, $4\times2$ table, one row per hub ($n=1{,}115$) \\
Smallest difference that matters & 5 percentage points in a format's gap rate \\
\bottomrule
\end{tabularx}}

$\pi_f$ is the share of format-$f$ hubs with no rows in 2014-H2 (overall 16.1\%). $\delta=5$ points, because a shift
that size changes the format mix of the hubs observed in 2014-H2 enough to bias any format comparison over that
period. Rule fixed in advance: the smallest expected count is 2.74 ($<5$, format~2 has 17 hubs), so the $p$-value comes
from 100{,}000 label permutations with the margins fixed; as a check, format~2 is merged into format~1.

\tpart{b}{Other tests we considered}
{\footnotesize
\begin{tabularx}{\linewidth}{@{}L{2.4cm} Y@{}}
\toprule
\textbf{Also considered} & \textbf{Why we did not use it} \\
\midrule
Two-proportion $z$ & tests one contrast, not four groups; used below for the follow-up estimate \\
Asymptotic $\chi^2$ only & an expected count below 5 weakens the approximation (we report it too) \\
Fisher's exact & standard form is $2\times2$; the permutation test is its Monte Carlo version for $4\times2$ \\
$\chi^2$ goodness of fit & compares one variable with fixed shares: the wrong question \\
\bottomrule
\end{tabularx}}

\tpart{c}{Result}
\[\chi^2=\sum_{\text{cells}}\frac{(O-E)^2}{E}=21.42,\]
where $O$ and $E$ are the observed and expected hub counts ($E=$ row total $\times$ column total$/n$), with
$(4-1)(2-1)=3$ d.f.\ and critical value $\chi^2_{0.95,3}=7.81$. Effect size: Cram\'er's
$V=\sqrt{\chi^2/(n\cdot\min(4-1,2-1))}=\sqrt{21.42/1115}=0.139$.
Gap rates are 14.0\%, 5.9\%, 9.5\% and 23.3\% for formats 1--4.
{\small
\begin{tabularx}{\linewidth}{@{}K{3.4cm} Y@{}}
\toprule
Estimate & format 4 vs rest: $\hat\pi_4-\hat\pi_{\text{rest}}=0.233-0.129=0.104$; Cram\'er's $V=0.139$ \\
95\% confidence interval & $0.053$ to $0.154$ (Wald, $SE=0.0257$; Newcombe $0.055$--$0.156$); $V$: $0.086$--$0.202$ \\
Test statistic (d.f.) & $\chi^2=21.42$ (3) \\
p-value & $2.1\times10^{-4}$ (permutation); $8.6\times10^{-5}$ (asymptotic) \\
Decision & \key{reject $H_0$} (also under Bonferroni) \\
\bottomrule
\end{tabularx}}

The whole interval lies beyond 5 points, so the imbalance is large enough to matter.
\textbf{Check and replication:} merged $3\times2$ table, $\chi^2(2)=20.62$, $p<10^{-4}$; odd and even HubID,
$\chi^2(3)=10.88$ ($p=0.015$) and $11.07$ ($p=0.012$).
\textbf{Implication:} the missingness is not random. Format~4 falls from 31.2\% of all hubs to 28.6\% of the hubs
observed in 2014-H2, so any comparison over that period leans towards formats 1--3. This is why we never impute the gap,
and why T4 requires 30 observed days in each window.

%------------------------------------------------------------ Test 3
\subsection*{Test 3 (P3): is competitor distance associated with typical volume?}

\tpart{a}{Hypotheses, written before running the test}
{\small
\begin{tabularx}{\linewidth}{@{}K{3.4cm} Y@{}}
\toprule
$H_0$ & $\rho=0$ \\
$H_1$ & $\rho\neq0$ \\
Significance level / side & $\alpha=0.05$, two-sided \\
Test & $t$-test on Pearson's $r$, one row per hub ($n=1{,}112$) \\
Smallest difference that matters & $|\rho|=0.10$ \\
\bottomrule
\end{tabularx}}

$\rho$ is the population correlation between $x_i=\ln(\text{distance})$ and $y_i=$ hub $i$'s mean $\log(1+V)$. Three
hubs with no distance are left out (complete case). Distance has skewness 2.93, so we use its log. $\delta=0.10$, because
an association explaining under 1\% of the variation between hubs would not inform where to site a hub.

\tpart{b}{Other tests we considered}
{\footnotesize
\begin{tabularx}{\linewidth}{@{}L{2.4cm} Y@{}}
\toprule
\textbf{Also considered} & \textbf{Why we did not use it} \\
\midrule
Spearman $\rho_s$ & robustness check (detects any monotone trend): $\rho_s=-0.037$, $p=0.21$ \\
Pearson on raw scale & a few far hubs dominate (skewness 2.93): $r=-0.045$, $p=0.14$ \\
Regression slope $t$ & gives the identical $t$ \\
Welch $t$, near vs far & a cut-off throws away information and is arbitrary \\
\bottomrule
\end{tabularx}}

\tpart{c}{Result}
\[t=\frac{(r-0)\sqrt{n-2}}{\sqrt{1-r^2}}=\frac{-0.08187\sqrt{1110}}{\sqrt{1-0.00670}}=-2.74.\]
{\small
\begin{tabularx}{\linewidth}{@{}K{3.4cm} Y@{}}
\toprule
Estimate & $r=-0.082$ (closer competitor, higher volume) \\
95\% confidence interval & $-0.140$ to $-0.023$ (Fisher $z$) \\
Test statistic (d.f.) & $t=-2.74$ (1{,}110); $t_{0.975,1110}=1.962$ \\
p-value & $0.0063$ \\
Decision & \key{reject $H_0$} (also under Bonferroni, $0.0063<0.0125$) \\
\bottomrule
\end{tabularx}}

\warn{The interval straddles $-0.10$, so we cannot tell whether the association is large enough to matter}; at most it
explains about 2\% of the variation. Controlling for format leaves it unchanged (partial $r=-0.082$). It replicates:
$r=-0.100$ on 2013 data and $-0.067$ on 2014--15 data. An earlier version of this test, Spearman on raw volume, failed to
reject ($\rho_s=-0.032$, $p=0.29$). The difference is explained by Figure~4 (a post-hoc look, not a test): the effect
comes from the closest decile alone (122 hubs at distance $\le250$, $+15.9\%$), while deciles 2--10 stay within $\pm4.4\%$.
With 5\% of each tail trimmed, $r=-0.055$ ($p=0.080$). A positive link between nearby competitors and demand is not
competition helping: competitors cluster where shoppers already are (dense shopping streets), a hidden third factor.

%------------------------------------------------------------ Test 4
\subsection*{Test 4 (P4): does a competitor opening change volume? (difference-in-differences)}

\tpart{a}{Hypotheses, written before running the test}
{\small
\begin{tabularx}{\linewidth}{@{}K{3.4cm} Y@{}}
\toprule
$H_0$ & $\mu_{\text{DiD}}=0$ \\
$H_1$ & $\mu_{\text{DiD}}\neq0$ \\
Significance level / side & $\alpha=0.05$, two-sided \\
Test & paired $t$-test on $\text{DiD}_i$, $n=144$ hubs \\
Smallest difference that matters & 3\% of volume (0.03 on the log scale) \\
\bottomrule
\end{tabularx}}

Only the opening \emph{month} is known, so that month is skipped. For a treated hub, $\Delta_i=\bar y_{\text{post}}-\bar
y_{\text{pre}}$ over the 56 days before and the 56 days after the opening month (at least 30 observed days in each).
$C_i$ is the same change averaged over control hubs in the same calendar windows, and $\text{DiD}_i=\Delta_i-C_i$, so
season, holidays and the shared promotion calendar cancel. Treated: openings from March 2013 to March 2015 with full
windows, 144 of the 180 openings inside the data. Controls: 498 hubs whose competitor opened before 2012 or after June
2015 (418--498 per window). The 1900 sentinel and the 354 unknown dates are left out. $\delta=3\%$, because a smaller fall
is within the hub's ordinary week-to-week variation and would not call for a response.

\tpart{b}{Other tests we considered}
{\footnotesize
\begin{tabularx}{\linewidth}{@{}L{2.4cm} Y@{}}
\toprule
\textbf{Also considered} & \textbf{Why we did not use it} \\
\midrule
Before/after, no control (our earlier design) & season is confounded with ``after'': the control change alone varies
  across windows with SD 0.093, as large as the effect. It gives $-5.6\%$ (CI $-7.8$ to $-3.4\%$) \\
Welch $t$, treated vs controls & ignores the pairing of each hub with its own calendar window \\
Wilcoxon signed-rank on $\text{DiD}_i$ & robustness check (the Q-Q plot shows a heavy left tail): $p=1.9\times10^{-17}$ \\
\bottomrule
\end{tabularx}}

\tpart{c}{Result}
\[t=\frac{\overline{\text{DiD}}-0}{s/\sqrt n}=\frac{-0.09218}{0.1170/\sqrt{144}}=\frac{-0.09218}{0.009752}=-9.45.\]
{\small
\begin{tabularx}{\linewidth}{@{}K{3.4cm} Y@{}}
\toprule
Estimate & $\overline{\text{DiD}}=-0.092$: \key{volume $8.8\%$ lower after an opening} \\
95\% confidence interval & $-0.111$ to $-0.073$ ($-10.5\%$ to $-7.0\%$) \\
Test statistic (d.f.) & $t=-9.45$ (143); $t_{0.975,143}=1.977$ \\
p-value & $9.0\times10^{-17}$ \\
Decision & \key{reject $H_0$} (also under Bonferroni) \\
\bottomrule
\end{tabularx}}

The whole interval lies beyond $-3\%$, so the fall is large enough to matter.
\textbf{Sensitivity and replication:} 28-day windows $-9.1\%$; 84-day windows $-8.4\%$; same-format controls $-8.8\%$;
reopening days excluded $-8.4\%$; odd and even HubID $-9.0\%$ and $-8.6\%$; openings in 2013 and in 2014--15 $-8.3\%$ and
$-9.3\%$. All reject.
\textbf{Placebo:} moving every opening 12 months earlier (92 hubs) gives $-0.7\%$ (CI $-2.0$ to $+0.6\%$), $t(91)=-1.08$,
$p=0.28$. We fail to reject, and the CI lies inside $\pm3\%$: the design does not invent a fall where none should exist.

\begin{figure}[!ht]
  \figfile{\FIG fig_qq.pdf}{3.0cm}
  \caption{T1 differences are close to normal; T4 differences have a heavy left tail, so we also report Wilcoxon (same
  decision).}
\end{figure}

\minihead{All four tests together}
{\footnotesize
\begin{tabularx}{\linewidth}{@{}L{0.55cm} L{1.75cm} L{1.4cm} Y L{1.25cm}@{}}
\toprule
\textbf{Test} & \textbf{Statistic (d.f.)} & \textbf{p} & \textbf{Estimate [95\% CI]} & \textbf{Matters?} \\
\midrule
T1 & $t=90.47$ (1{,}114) & $<10^{-300}$ & $+30.9\%$ [30.1, 31.6] (ratio of multipliers) & \key{yes} \\
T2 & $\chi^2=21.42$ (3) & $2.1\times10^{-4}$ & $+10.4$ points [5.3, 15.4] & \key{yes} \\
T3 & $t=-2.74$ (1{,}110) & $0.0063$ & $r=-0.082$ [$-0.140$, $-0.023$] & \warn{uncertain} \\
T4 & $t=-9.45$ (143) & $9.0\times10^{-17}$ & $-8.8\%$ [$-10.5$, $-7.0$] & \key{yes} \\
\bottomrule
\end{tabularx}}

All four raw $p$-values are below the Bonferroni level 0.0125, and Holm gives the same four rejections. \warn{T3 is the fragile one: it rests on one decile.}
\key{What would change our mind:} T1, a week-level interval reaching below 0.05; T3, the closest-decile effect
vanishing once a measure of shopper density is added; T4, placebo dates producing falls of similar size.

%-----------------------------------------------------------------------
\section{Limitations}

\begin{itemize}
  \item \key{Association, not cause.} The business sets one promotion calendar for all hubs, so promotion weeks may
  differ from other weeks in ways we cannot see. Competitors choose where to open, and T3 suggests they choose busy
  places. DiD assumes treated and control hubs would have moved in parallel; the placebo supports this but cannot prove it.
  \item \key{Shared dates.} All hubs share the calendar, so hub-level standard errors ignore cross-hub correlation.
  The week-level T1 check had a standard error 9 times larger. T4's control subtraction removes shared date shocks;
  T2 and T3 use one static value per hub.
  \item \key{Rows left out.} 166{,}269 closed and 54 open-zero rows; the non-random 2014-H2 gap (T2); 36 of the 180
  openings (incomplete windows) and the 354 hubs with an unknown opening date (T4); 3 hubs without a distance (T3).
  \item \key{Month-precision event dates.} The opening month is skipped, so a fall that starts and ends within it
  would be missed.
  \item \key{Provenance and units.} The data is renamed public data, and the unit of \code{OrderVolume} is unstated.
  Our effects are relative (\%), so they do not depend on the unit.
  \item \key{Choices made after EDA.} The questions came after exploration and an earlier raw-scale run, so all tests
  are two-sided and the plan is not a true pre-registration. The T3 decile breakdown is post hoc.
  \item \key{Scope.} January 2013 to June 2015 only; the test file's summer weeks are covered only through earlier
  summers.
\end{itemize}

\end{document}
